% !TeX program = xelatex
% Сборка: xelatex Dmitrich.tex
\documentclass[11pt,letterpaper]{article}

\usepackage[letterpaper,margin=0.5in,bottom=0.42in]{geometry}
\usepackage{fontspec}
\setmainfont{NewCM10-Regular.otf}[
  BoldFont       = NewCM10-Bold.otf,
  ItalicFont     = NewCM10-Italic.otf,
  BoldItalicFont = NewCM10-BoldItalic.otf
]
\usepackage{babel}
\babelprovide[main,import]{russian}
\usepackage{enumitem}
\usepackage[unicode,hidelinks]{hyperref}
\hypersetup{
  pdftitle  = {Медведев Даниил — резюме},
  pdfauthor = {Даниил Медведев},
  pdflang   = {ru-RU}
}

\pagestyle{empty}
\setlength{\parindent}{0pt}
\setlength{\parskip}{0pt}
\raggedright
\setlist[itemize]{
  leftmargin=0.50in, labelsep=0.09in, rightmargin=0.09in,
  label={\raisebox{0.08ex}{\scriptsize\textbullet}},
  topsep=2pt, itemsep=1.1pt, parsep=0pt, partopsep=0pt
}

\newcommand{\cvsection}[1]{%
  \par\addvspace{10pt}%
  {\Large\scshape #1\par}\nobreak
  \vspace{4.8pt}\hrule height 0.4pt\nobreak\vspace{4pt}%
}
\newcommand{\entry}[2]{%
  \par\addvspace{3pt}\hspace*{0.15in}\textbf{#1}\hfill #2\hspace*{0.08in}\par
}
\newcommand{\detail}[2]{%
  \hspace*{0.15in}\textit{#1}\hfill\textit{#2}\hspace*{0.08in}\par
}
\newcommand{\linktext}[2]{\href{#1}{\underline{#2}}}

\begin{document}

{\centering
  {\fontsize{25}{29}\selectfont\scshape Медведев Даниил\par}
  \vspace{8pt}
  \linktext{mailto:yalkar2006@mail.ru}{yalkar2006@mail.ru}
  \textbar{}
  \linktext{https://t.me/yalkar}{https://t.me/yalkar}
  \textbar{}
  \linktext{https://github.com/yalkar}{https://github.com/yalkar}\par
}
\vspace{4pt}

\cvsection{Образование}
\entry{Высшая школа экономики (НИУ ВШЭ, МИЭМ)}{Россия, Москва}
\detail{Информационная безопасность, бакалавриат, 4-й курс}{2023--2027}
\begin{itemize}
  \item \textbf{Релевантные курсы} — администрирование систем и сетей, криптографические
    методы защиты информации, разработка защищённых приложений и протоколы безопасности,
    основы криптографии, алгоритмизация и программирование, администрирование ОС Astra Linux,
    системы и сети передачи информации.
\end{itemize}

\cvsection{Опыт работы}
\entry{Специалист по информационной безопасности}{По настоящее время}
\detail{Банковский сектор}{Безопасность партнёрских подключений}
\begin{itemize}
  \item Собрал на Python (pandas) витрину партнёрских VPN-подключений. Сопоставил
    YAML-реестры из GitLab со сведениями о договорах и правилами межсетевых экранов.
  \item Настроил оценку риска туннелей по направлению трафика, передаче персональных
    данных и доступу к критичным ресурсам. Добавил проверку признаков файлового обмена.
  \item Создал BI-дашборд с ежедневной историей. По метрикам на SQL команда отслеживает
    просроченные доступы и подключения без ссылки на договор.
  \item Подключил данные NDR для выявления неиспользуемых доступов и настроил в n8n
    создание задач на их отзыв в Jira.
  \item Автоматизировал создание merge request для обновления реестра в GitLab.
\end{itemize}

\cvsection{Проекты}
\entry{Изучение возможности размещения скрытых уникальных меток в печатных изображениях}
  {\small\linktext{https://github.com/yalkar/watermark-method-stage8}{Код и результаты}\quad\normalsize 2026}
\begin{itemize}
  \item Реализовал на Python метод Thongkor (блочное DCT) для встраивания и извлечения
    водяных знаков. Адаптировал метод на основе преобразования Фурье (DFT) для многобитных меток.
  \item На 100 изображениях проверил влияние JPEG-сжатия, шума и поворота. На 150 снимках
    50 карточек измерил долю битовых ошибок (BER) при извлечении 256-битной метки после печати.
\end{itemize}

\entry{Система обработки событий ИБ}
  {\small\linktext{https://github.com/yalkar/security-event-processing-system}{Код на GitHub}\quad\normalsize 2026}
\begin{itemize}
  \item Разработал сервис на FastAPI с очередью RabbitMQ, хранилищем PostgreSQL и кэшем
    Redis. Разделил контейнеры на четыре сети и проверил восстановление базы из резервной копии.
\end{itemize}

\entry{Исследование безопасности в IoT-устройствах}{2025}
\begin{itemize}
  \item Систематизировал основные угрозы безопасности IoT и способы их предотвращения.
  \item Изучил сетевые сервисы и трафик с помощью Nmap и Wireshark, подготовил
    рекомендации по защите IoT-устройств.
\end{itemize}

\cvsection{Технические навыки}
\begin{list}{}{\setlength{\leftmargin}{0.15in}\setlength{\rightmargin}{0.08in}%
  \setlength{\topsep}{0pt}\setlength{\itemsep}{1pt}\setlength{\parsep}{0pt}}
  \item \textbf{Языки программирования и запросов} — Python, C++, SQL.
  \item \textbf{Рабочие инструменты} — pandas, Git, GitLab API, YAML, n8n, Jira,
    VPN/IPsec, правила межсетевых экранов, NDR, NetBox (pynetbox).
  \item \textbf{Сети и защита (учебная практика)} — TCP/IP, ACL, Wireshark, Nmap,
    Suricata (IDS/IPS), Sigma, Metasploit.
  \item \textbf{Инфраструктура и разработка (проекты)} — Linux (Astra Linux), Docker Compose,
    GitLab CI/CD, Kubernetes (NetworkPolicy), FastAPI, REST API, SOAP, PostgreSQL, Redis,
    RabbitMQ, NumPy, OpenCV.
\end{list}

\end{document}
